% Part: proof-theory
% Chapter: natural-deduction
% Section: introduction

\documentclass[../../../include/open-logic-section]{subfiles}

\begin{document}

\olfileid{pt}{nat}{int}

\olsection{Pambuka}

Deduksi alami iku kulawarga sistem !!{proof} kang saben
konektif logis utawa kuantor nduwèni pasangan inferènsi:
siji ``nglebokaké'' operator mau (aturan
!!a{introduction}) lan siji ``ngilangaké'' operator
mau (aturan !!a{elimination}). Tuladhané,
\Intro{\lif} nduwèni !!{formula}s kanthi wujud
~$!A \lif !B$ minangka konklusi, déné \Elim{\lif}
nduwèni !!{formula}s kanthi wujud~$!A \lif !B$
minangka premis.

Rong versi deduksi alami kang sepisanan dikenalaké ing
taun 1930-an dening Gerhard Gentzen lan Stanisław Jaśkowski.
Sistem Gentzen kang paling akèh dirembug ahli teori bukti,
dene saka Jaśkowski lair sistem kang paling akèh dienggo
mulang. Ing kaloro sistem, bukti bisa diwiwiti saka
\emph{asumsi}---!!{formula}s kang ora perlu
!!{prove}d, nanging bisa dienggo kanggo !!{proving}
!!{formula}s liya. Sakabèhé !!{proof} bisa gumantung
marang asumsi mau, lan amarga asumsi durung kabuktèkaké,
!!{proof} iku ora langsung netepaké konklusiné
(\emph{!!{formula} pungkasan}), nanging mung nuduhaké
yèn konklusiné ndhèrèk saka asumsi.

Sawetara aturan inferènsi nduwèni konklusi kang ora
gumantung manèh marang sawetara asumsi: aturan mau
``mbusak'' asumsi. Tuladhané, aturan \Intro{\lif}
njupuk !!a{proof} kanggo konklusi~$!B$ saka
asumsi~$!A$, banjur ngasilaké $!A \lif !B$ minangka
konklusi. !!{proof} kang pungkasané ana ing
$!A \lif !B$ ora gumantung manèh marang~$!A$;
asumsi $!A$ wis dibusak. Ing sistem Gentzen,
!!{proof}s wujudé wit !!{formula}s, asumsi ana minangka
!!{formula}s paling dhuwur, lan aturan kaya~\Intro{\lif}
nduwèni label kang nuduhaké asumsi endi kang dibusak.
Ing sistem Jaśkowski, pérangan !!{proof} kang gumantung
marang asumsi dipisah kanthi cara tartamtu, upamané
digeser mlebu ing mburi garis vertikal utawa dilebokaké
ing kothak, kanthi asumsi~$!A$ kang dibusak lan
konsekuèn $!B$ saka implikasi minangka larik kapisan
lan pungkasan ing kothak, déné konklusi~$!A \lif !B$
saka \Intro{\lif} ana ing njaban kothak.

Sistem iki diarani ``alami'' amarga aturan inferènsiné
niru carané kita (utawa paling ora matématikawan)
mikiran kanthi alami. Kanggo netepaké asil hipotètis
(implikasi), kita nganggep antesèdené bener lan
migunakaké kanggo mbuktèkaké konsekuèné.
Iki~\Intro{\lif}. Yèn kita ngerti implikasi
$!A \lif !B$ lan antesèdené~$!A$, kita nyimpulaké~$!B$.
Iki~\Elim{\lif}. Kita migunakaké bukti miturut kahanan
kanggo nuduhaké yèn saka disjungsi~$!A \lor !B$
ndhèrèk sawijining konklusi. Iki~\Elim{\lor}.
Kanggo mbuktèkaké pratelan umum~$\lforall[x][!A(x)]$,
kita mbuktèkaké $!A(c)$ kanggo~$c$ kang arbitrèr.
Iki~\Intro{\lforall}. Lan saterusé.

Tuladhané, iki !!{proof} prasaja kanggo
~$!A \lif (!A \lor !B)$ ing deduksi alami gaya Gentzen:
\[
\AxiomC{$\Discharge{!A}{x}$}
\RightLabel{\Intro{\lor}}
\UnaryInfC{$!A \lor !B$}
\DischargeRule{\Intro{\lif}}{x}
\UnaryInfC{$!A \lif (!A \lor !B)$}
\DisplayProof
\]
Bukti iki diwiwiti saka asumsi~$!A$, nganggo
\Intro{\lor} kanggo nyimpulaké~$!A \lor !B$,
banjur \Intro{\lif} kanggo nyimpulaké
~$!A \lif (!A \lor !B)$. Inferènsi \Intro{\lif}
mbusak asumsi~$!A$; iki ditandhani label~$x$.

Ahli teori bukti luwih seneng sistem asli Gentzen kang
nganggo wit !!{formula}s, nanging ana varian liya kang
uga wigati. !!^a{proof} kanggo~$!A$ saka
asumsi~$\Gamma$ nuduhaké yèn $!A$ ndhèrèk saka~$\Gamma$,
yaiku $\Gamma \Entails !A$. Aturan inferènsi deduksi
alami bisa diwaca minangka pratelan babagan akibat
mangkono; tuladhané, \Intro{\lif} mratélakaké yèn
$\Gamma + !A \Entails !B$, banjur
$\Gamma \Entails !A \lif !B$. Dadi wajar yèn aturan
dirumusaké langsung kanggo nggarap asumsi lan
konklusi kang ndhèrèk saka asumsi mau ing premis
lan konklusi aturan inferènsi, yaiku nggarap sekèn
$\Gamma \Sequent !A$. Bukti prasaja ing ndhuwur
dadi kaya mangkéné ing sistem iku:
\[
\Axiom$x:!A \fCenter !A$
\RightLabel{\Intro{\lor}}
\UnaryInf$x:!A \fCenter !A \lor !B$
\DischargeRule{\Intro{\lif}}{x}
\UnaryInf$\fCenter !A \lif (!A \lor !B)$
\DisplayProof
\]
Kaya kang katon, aturané padha: aturan nggarap
!!{formula} ing sisih tengen~$\Sequent$, déné
!!{formula}s ing sisih kiwa mung nyathet asumsi
kang dienggo ing saben langkah.

Pasangan aturan !!{introduction}/!!{elimination}
waé durung cukup kanggo logika klasik. Kita perlu
nambah rong aturan kang gegayutan karo~$\lfalse$.
Kapisan yaiku ``aturan absurditas intuisionistik''~\FalseInt,
utawa ``èksplosi'', kang ngidini nyimpulaké apa waé
saka $\lfalse$. Kapindho yaiku prinsip klasik bukti
ora langsung~\FalseCl{} (utawa ``aturan absurditas
klasik''), kang mratélakaké yèn kita bisa
!!{prove}~$\lfalse$ saka~$\lnot !A$, kita wis
!!{prove}~$!A$. Yèn~\FalseCl dibusak, kita ngolèhké
sistem kang cocog kanggo logika intuisionistik.
Yèn kaloroné dibusak, kita ngolèhké kang diarani
``logika minimal''.

Kita bakal ngrembug deduksi alami gaya Gentzen kanggo
logika klasik lan intuisionistik. Iki sistem \Log{N1c}
lan \Log{N1i}. Sistem \Log{N2c} lan \Log{N2i} iku
versi kang cocog nanging nggarap sekèn
$\Gamma \Sequent !A$ tinimbang !!{formula}s siji-siji~$!A$.

\end{document}
